\documentclass{article}
\usepackage{graphicx, physics, dsfont}
\usepackage{fancyhdr}
\usepackage{hyperref}
\usepackage{ragged2e}
\usepackage{amsmath, amsthm, amssymb}
\usepackage[margin=1in]{geometry}
\usepackage{setspace}
\usepackage{tikz}
\usetikzlibrary{positioning}

\pagestyle{fancy}
\lhead{Zoeb Izzi}
\rhead{\today}
\lfoot{Problem Set 9}
\rfoot{}

\begin{document}
\thispagestyle{fancy}
\begin{center}\LARGE{Advanced Mathematical Techniques \\ Problem Set 9}\end{center}
\vspace{10mm}
\begin{enumerate}
    \item In a very strange gambling hall, there is a game in which gamblers can place bets on how many 1's will come up in a run of 24 rolls of a fair six-sided die.
    \begin{enumerate}
        \item What is the expectation value of the number of times a 1 will come up?
        \\
        \[24 \times \frac 1 6 = \boxed 4\]
        \item What is the probability that a 1 will come up its expected number of times?
        \\
        \[\binom {24}{4} \times \left(\frac 1 6\right)^4 \left( \frac 5 6 \right)^{20} \approx \boxed{0.21386}\]
        \\Note that the following parts (c)-(g) refer to a gambler who places a bet that 1 will come up four times \textit{or more}. 
        \item At the beginning of the run of rolls, what is the probability that the gambler will win? 
        \\
        \[1 - \left( \frac 5 6\right)^{24} - \binom{24}{1}\left(\frac 1 6\right) \left(\frac 5 6\right)^{23} - \binom{24}{2}\left(\frac 1 6\right)^2 \left(\frac 5 6\right)^{22} - \binom{24}{3}\left(\frac 1 6\right)^3 \left(\frac 5 6\right)^{21} \approx \boxed{0.5845}\]
        \item Suppose that ten rolls have already taken place and 1 has come up twice. What is the probability that the gambler will win given this information?
        \\
        \[1 - \left( \frac 5 6\right)^{14} - \binom{14}{1}\left(\frac 1 6\right) \left(\frac 5 6\right)^{13}  \approx \boxed{0.7040}\]
        \item Suppose that ten rolls have already taken place and 1 has come up three times. What is the probability that the gambler will win given this information?
        \\
        \[1 - \left( \frac 5 6\right)^{14} \approx \boxed{0.9221}\]
        \item Suppose that ten rolls have already taken place and none of them resulted in a 1. What is the probability that the gambler will win given this information?
        \\
        \[1 - \left( \frac 5 6\right)^{14} - \binom{14}{1}\left(\frac 1 6\right) \left(\frac 5 6\right)^{13} - \binom{14}{2}\left(\frac 1 6\right)^2 \left(\frac 5 6\right)^{12} - \binom{14}{3}\left(\frac 1 6\right)^3 \left(\frac 5 6\right)^{11} \approx \boxed{0.1937} \]
        \item Explain how the additional information in parts (d)-(f) changes the likelihood of the gambler being successful.
        \\ \\ Changing the amount of rolls that have already occurred reduces the probability by reducing the amount of possible tries the gambler has to roll the 1s needed to win, and increasing the number of 1s already rolled reduces the amount that the gambler needs rolled in order to win. \\
        \item Suppose now that the die will be rolled 20 times instead of 24. How does this change your answers to (a) and (b)? What does this imply about the statement, "The largest probability is associated with 1 coming up the expected number of times."?
        \\ \\ Changing the amount of trials to be done makes my answer to part (a) $20 \times \frac 1 6 = \boxed{\frac {10} 3}$ and part (b) $\binom{20}{4} \left( \frac 1 6\right)^4 \left(\frac 5 6 \right)^{16} \approx \boxed{0.2022}$. This implies that the statement is not necessarily true, since 1 can't possibly come up a fractional number of times.
    \end{enumerate}
    \item You measure a certain length 4 times and obtain the results 3.26, 3.28, 3.26, and 3.27 cm. Your measurement device indicates a precision of 0.005 cm, but you want to use your repeated measurements to gauge the \textit{actual} uncertainty in this length.
    \begin{enumerate}
        \item Use this information to estimate the expectation value and the standard deviation of the length.
        \\
        \[\langle k \rangle = \frac{3.26 + 3.28 + 3.26 + 3.27}{4} = \boxed{3.2675}, \sigma_k = \sqrt{\frac{1}{3}\sum_{i=1}^4 (x_i - \langle k \rangle)^2} \approx \boxed{0.009574}\]
        \item Determine the standard devation of the mean associated with this length and quote your resulting length as "number $\pm$ error".
        \\
        \[\text{number $\pm$ error} = \boxed{3.2675 \pm 0.004787}\]
        \item Use the normal distribution to estimate the probability that a fifth measurement of this length results in any value $\le 3.25$ cm.
        \\
        \[z = \frac{3.25-3.2675}{0.009574} \approx -1.8278, P(X \le 3.25) \approx \boxed{0.03379}\]
        \item Suppose a 5th measurement was taken before the results were analyzed and the result was 3.25 cm. Re-do parts (a)-(c) with this new data and comment on any differences seen in the new part (c). How is this situation similar to that found in part (f) on problem 1?
        \\ 
        \[\langle k \rangle = \frac{3.26 + 3.28 + 3.26 + 3.27 + 3.25}{5} = \boxed{3.264}, \sigma_k = \sqrt{\frac{1}{4}\sum_{i=1}^5 (x_i - \langle k \rangle)^2} \approx \boxed{0.01140}\]
        \[\text{number $\pm$ error} = \boxed{3.264 \pm 0.005701}\]
        \[z = \frac{3.25-3.264}{0.01140} \approx -1.2279, P(X \le 3.25) \approx \boxed{0.1097}\]
        This is similar to the situation in part (f) of problem 1 because the additional information of the 5th measurement being 3.25 cm increases the standard deviation and lowers the mean, which makes it more likely that a random measurement will result in 3.25 cm or less.
        \end{enumerate}
    \item A room contains 7 people. You are interested in determining the probability of various distributions of the day of the week on which the people were born. Assume all days of the week are equally likely.
    \begin{enumerate}
        \item Determine the probability that all seven people were born on different days of the week.
        \\ 
        \[  \frac{7!}{7!} \cdot \binom{7}{1,1,1,1,1,1,1} \cdot \frac 1 {7^7} = \frac{7!}{7^7} \approx \boxed{0.006120}\]
        \item Determine the probability that exactly two people were born on the same day of the week and everyone else was born on different days.
        \\
        \[\frac{7!}{1!5!1!} \cdot \binom{7}{2,1,1,1,1,1} \cdot \frac{1}{7^7} = \frac{42 \cdot 7!}{2 \cdot 7^7} = \boxed{0.1285}\]
        \item Determine the probability that exactly two \textit{pairs} of people were born on the same day of the week and everyone else was born on different days.
        \\
        \[\frac{7!}{2!2!3!} \cdot \binom{7}{2,2,1,1,1} \cdot \frac{1}{7^7} = \frac{210 \cdot 7!}{4\cdot 7^7} \approx \boxed{0.3213} \]
        \item Determine the probability that exactly three people were born on the same day of the week and all of the other people were born on different days.
        \\
        \[\frac{7!}{1!4!2!} \cdot \binom{7}{3,1,1,1,1} \cdot \frac{1}{7^7} = \frac{35 \cdot 7!}{5 \cdot 7^7} \approx \boxed{0.1071}\]
        \item Show that the sum of the results for parts (a)-(d) is less than 1 and indicate what distributions the remaining probability consists of. Pick one of these and show that its probability is smaller than the difference between your sum and 1.
        \\
        \\ The remaining probability consists of other partitions of 7 people, which include one triple, one pair, and two singles, three pairs and one single, one quadruple and three singles (indicated below), two triples and one single, etc.
        \[0.006120 + 0.1285 + 0.3213 + 0.1071 \approx 0.5630 < 1\]
        \[P(\text{4 born same day, all else different days}) = \frac{7!}{1!3!3!}\cdot\binom{7}{4,1,1,1}\cdot\frac{1}{7^{7}} \approx \boxed{0.03570} < 1 - 0.5630 = 0.4370\]
        \item Determine the expectation value of the product of the number of people born on Tuesday and the number of people born on Saturday. Why isn't this expectation value just equal to 1? Explain.
        \\\\This expectation value isn't just 1 because when someone is born on Tuesday, that leaves one less person who can possibly be born on Saturday, so the expectation value here is less than 1.
    \end{enumerate}
\end{enumerate}
\end{document}